\documentclass[11pt]{article}
\usepackage[a4paper,margin=25mm]{geometry}
\usepackage{amsmath,amssymb,amsthm,mathtools,microtype,titling}
\usepackage[hidelinks]{hyperref}
\setlength{\droptitle}{-8mm}
\newtheorem{theorem}{Theorem}
\newtheorem{lemma}{Lemma}
\newtheorem{corollary}{Corollary}
\newcommand{\E}{\mathbb E}
\newcommand{\R}{\mathbb R}
\newcommand{\eps}{\varepsilon}
\newcommand{\cA}{\mathcal A}
\newcommand{\cD}{\mathcal D}
\DeclareMathOperator{\Var}{Var}
\DeclareMathOperator{\Cov}{Cov}
\title{Shock selection in a nonpotential\protect\\ finite-state master equation}
\author{Ji Ho Bae\\\small LawAll AI}
\date{10 October 2026}
\hypersetup{pdftitle={Shock selection in a nonpotential finite-state master equation},pdfauthor={Ji Ho Bae}}
\begin{document}
\maketitle
\begin{abstract}
We give an explicit family of smooth nonpotential costs for which the uncorrected finite-state master equation has a unique selected weak feedback field. Every admissible vanishing boundary repulsion gives the same local $L^1$ limit, without a rate condition between the repulsion and the common noise. The limit has an intrinsic scalar entropy characterization. The family includes interior shocks: its permitted time horizon is independent of front steepness, and a scaling extends the result to every finite horizon. The proof uses localization of the passive values and a Cole--Hopf comparison construction that remains valid through shock formation.
\end{abstract}

\section{The problem and a class of data}
Fix $d\ge3$ and the simplex $\Delta_d=\{m_i\ge0:\sum_i m_i=1\}$. Write
\[
 H_i(U)=\frac12\sum_{j\ne i}(U_i-U_j)_+^2,\qquad
 b_i(m,U)=\sum_{j\ne i}\{m_j(U_j-U_i)_+-m_i(U_i-U_j)_+\}.
\]
The uncorrected regularization considered here is
\begin{align}\label{master}
0={}&\partial_tU_i+\eps\sum_{j,l}(m_j\delta_{jl}-m_jm_l)\partial^2_{m_jm_l}U_i
 +(b+R^\phi)\cdot DU_i+2\eps(e_i-m)\cdot DU_i-H_i(U)\notag\\
&+f_i(m)+\sum_{j\ne i}\phi_\eps(m_j)(U_j-U_i),\qquad U_i(T,m)=g_i(m),
\end{align}
where $R_i^\phi=\phi_\eps(m_i)-m_i\sum_j\phi_\eps(m_j)$ and derivatives are tangent to the simplex. Following the selection formulation in catalogue C73-30, let $\cA(f,g,T)$ consist of smooth nonnegative families $\phi_\eps$, converging to zero uniformly on every compact subset of $(0,1]$, for which \eqref{master} has a unique classical solution continuous on the closed cylinder, with bounded value and first derivatives and locally bounded second derivatives in the interior. No bound on these derivatives uniform in $\eps$ is imposed. The feedback field is $(U_i-U_d)_{i<d}$.

The problem asks for nonpotential assumptions and an intrinsic admissibility condition giving existence of selected limits, uniqueness independent of the admissible regularization, and agreement with classical solutions. The potential theory and this extension are discussed in \cite[\S4.3--4.4]{CD}. We address the sufficient-condition question with a family that undergoes loss of classical regularity.

Put $q=(d-1)/2$, $p=m_1$, and $r=(r_2,\ldots,r_{d-1})$, where $r_j=m_j-m_d$. Thus
\[
 m_d=\frac{1-p-\sum_{j=2}^{d-1}r_j}{d-1},\qquad m_j=m_d+r_j.
\]
Choose fixed smooth cutoffs $\chi:\R\to[0,1]$ and $\beta:\R^{d-2}\to[0,1]$ such that
\[
 \chi=0\ \text{on }(-\infty,1/128],\quad \chi=1\ \text{on }[1/64,\infty),
\]
and, with $R_d=1/(16d)$, $\beta=1$ on $\{|r|_\infty\le R_d/4\}$ and $\beta=0$ on $\{|r|_\infty\ge R_d/2\}$. Set $\alpha=1/16$, $y=\sqrt p$, $\theta=y-1/2-\alpha r_2$, and, for $k\ge1$, define
\begin{equation}\label{data}
 f=0,\qquad g=(G_k,0,\ldots,0),\qquad
 G_k(p,r)=\frac1q\left[24(1-\chi(p))+\chi(p)\frac{2-\beta(r)\tanh(k\theta)}{\sqrt p}\right].
\end{equation}
The second term is defined as zero where $\chi=0$. These data are smooth on a neighborhood of $\Delta_d$, and $1/q\le G_k\le34/q$ there on the simplex.

For each admissible $r$, let $0<y<L(r)$ be the interval corresponding to the simplex interior, where
\[
 L(r)^2=1-\sum_{j=2}^{d-1}r_j+(d-1)\min(0,r_2,\ldots,r_{d-1})>0.
\]
Define $v_0(y,r)=qyG_k(y^2,r)$, extend it by zero to $y<0$, and extend it constantly to the right of $L(r)$. The construction gives $0\le v_0\le3$. Let $v(s,y,r)$ be the bounded nonnegative entropy solution of
\begin{equation}\label{burgers}
 v_s+\partial_y(v^2/2)=0,\qquad v(0,y,r)=v_0(y,r).
\end{equation}
The value on $0<y<L(r)$ is independent of the chosen nonnegative right extension. Indeed, its Hopf--Lax minimizers lie in $[0,y]$: the initial primitive is nondecreasing, and the quadratic penalty increases for a footpoint to the right of $y$. The zero extension also gives the left trace zero. We use the usual bounded entropy uniqueness and vanishing-viscosity theorem for this scalar Cauchy problem \cite{Kru}.

\begin{theorem}\label{main}
For each $d\ge3$ and the fixed cutoffs above, there is a $T_0>0$, independent of $k\ge1$, with the following properties. For $0<T\le T_0$ and the data \eqref{data}, the class $\cA(f,g,T)$ is nonempty. For every $\phi\in\cA(f,g,T)$,
\begin{equation}\label{convergence}
 U^{\eps,\phi}\longrightarrow
 U^*(t,m)=\left(\frac{v(T-t,\sqrt p,r)}{q\sqrt p},0,\ldots,0\right)
 \quad\text{in }L^1_{\rm loc}((0,T)\times\operatorname{int}\Delta_d).
\end{equation}
Consequently all selected feedback fields coincide, and any two admissible families have feedback distance tending to zero on every interior compact set.

Among bounded feedback fields $Z=(Z_1,\ldots,Z_{d-1})$, an intrinsic admissibility condition is that $Z_2=\cdots=Z_{d-1}=0$, $Z_1\ge0$, and that for almost every $r$ the function $q\sqrt p\,Z_1$, expressed in time $s=T-t$ and in $y=\sqrt p$, is an entropy solution of \eqref{burgers} on $0<y<L(r)$ with initial trace $v_0(\cdot,r)$: its $L^1$ distance to $v_0(\cdot,r)$ on each compact subinterval of $(0,L(r))$ tends to zero as $s\downarrow0$ outside a null set of times. This condition determines the selected field uniquely and agrees with the classical master field whenever the latter exists.

The data are nonpotential even modulo a common component. If $kT>1$, the selected field has an interior shock for $1/k<s\le T$. Thus the assertion includes a nonclassical regime, while retaining every family in $\cA$.
\end{theorem}

\section{Bounds that do not use derivatives uniform in the noise}
The following estimates also justify the reduction to one active value. All constants may depend on the fixed data and on the compact set, but not on $\eps$ or on derivatives of $\phi_\eps$.

\begin{lemma}\label{localization}
Suppose $f=0$, $g=(G,0,\ldots,0)$ and $0<\gamma\le G\le M$. Every admissible solution satisfies $0\le U_i\le M$. For $D_\eta=\{m_i\ge\eta\ \forall i\}$, put
\[
 C=(d-1)M+1,\quad \delta=\tfrac12\eta e^{-CT},\quad
 a_\eps=\sup_{[\delta,1]}\phi_\eps,\quad
 e_\eps=d\exp\left(-\frac{\eta^2}{4\eps T e^{2CT}}\right).
\]
For all sufficiently small $\eps$, on $[0,T]\times D_\eta$,
\begin{equation}\label{passive}
 0\le U_i\le M(Ta_\eps+e_\eps)\quad(i>1),\qquad
 U_1\ge\gamma e^{-[qM+(d-1)a_\eps]T}(1-e_\eps).
\end{equation}
\end{lemma}
\begin{proof}
First fix $\eps$. The ordinary maximum principle extends to the degenerate boundary as follows. Add $hB_a(m)=h\sum_jm_j^a$, $0<a<1$, to a component being maximized. Bounded first derivatives force the maximum into the interior: adding a small empty coordinate gains order $h\rho^a$, whereas the other changes are $O(\rho)$. At this maximum the Hamiltonian and value-switching terms have the appropriate nonpositive signs. The drift without $R^\phi$ has coordinate lower bounds $-C_\eps m_j$, since the solution is bounded for this fixed $\eps$. The repulsion contributes
\[
 -hR^\phi\cdot DB_a=-ha\sum_j\phi_\eps(m_j)\{m_j^{a-1}-B_a(m)\}.
\]
Because $B_a\le d^{1-a}$, each summand with $m_j\le1/d$ is nonpositive. The remaining coordinates lie in $[1/d,1]$, where $\phi_\eps$ is bounded. Thus this contribution is $O_\eps(h)$ as well, regardless of the size of the repulsion close to zero.

If $m_l\ge1/d$, the tangent first-order condition bounds
$ha m_j^{a-1}$ by $L_\eps+ha d^{1-a}$, where $L_\eps$ bounds the first derivatives. The diffusion error is at most
\[
 \eps d(1-a)(L_\eps+ha d^{1-a}),
\]
and the drift error is $O_\eps(h)$. A strict time barrier, followed by $h\downarrow0$ and then $a\uparrow1$, gives $U_i\le M$. Applied to $-U$ it gives $U_i\ge0$; at a minimum component $H_i=0$. No uniform bound on $L_\eps$ was used.

Stop an interior diffusion, with the Kimura covariance and drift
$b(m,U)+R^\phi+2\eps(e_I-m)$, on leaving $D_\delta$ or reaching $T$. The label $I$ may switch in any adapted way. Once $(d-1)a_\eps+2\eps\le1$, each drift coordinate is at least $-Cm_j$. After multiplication by $e^{C(s-t)}$, its finite-variation part is nondecreasing. Its martingale part has quadratic variation at most $\eps T e^{2CT}/2$, because the coordinate covariance rate is $2\eps m_j(1-m_j)\le\eps/2$. An exit requires one such martingale to fall by at least $\eta/2$. The exponential martingale inequality and a union bound give exit probability at most $e_\eps$.

For a passive component, write
\[
 -H_i(U)=\sum_{j\ne i}c_{ij}(U_j-U_i),\qquad c_{ij}=\tfrac12(U_i-U_j)_+.
\]
Its $j=1$ term is nonpositive and can be dropped for an upper comparison. The repulsion term involving $U_1$ is at most $M\phi_\eps(m_1)$. Keep the nonnegative switching rates between passive labels. The stopped cooperative-system representation has terminal value zero, exit value at most $M$, and source at most $Ma_\eps$. This proves the first estimate in \eqref{passive}. All coefficients are bounded on the stopped interior domain for fixed $\eps$.

Finally $H_1\le qU_1^2\le qMU_1$, and its repulsion value term is at least $-(d-1)a_\eps U_1$. The stopped scalar comparison with killing rate $qM+(d-1)a_\eps$, terminal value $\gamma$ and exit value zero gives the second estimate.
\end{proof}

We also record a local consistency consequence used at the boundary of the shock region. Write $A=(m_j\delta_{jl}-m_jm_l)_{j,l}$, so that the second-order part of \eqref{master} is $\eps\operatorname{tr}(AD^2U_i)$, and $|x|_1=\sum_l|x_l|$ for $x\in\R^d$.

\begin{lemma}\label{consistency}
Suppose $f=0$, $g=(G,0,\ldots,0)$, $0<\gamma\le G\le M$ and $\phi\in\cA(f,g,T)$. Let $\mathcal O\subset\operatorname{int}\Delta_d$ be open, and let $U^{\rm cl}$, once continuously differentiable in $t$ and twice in $m$, solve the deterministic master equation
\[
 \partial_tU^{\rm cl}_i+b(m,U^{\rm cl})\cdot DU^{\rm cl}_i-H_i(U^{\rm cl})=0,\qquad U^{\rm cl}(T)=g,
\]
on a neighborhood of $[0,T]\times(\operatorname{int}\Delta_d\setminus\mathcal O)$ in $[0,T]\times\operatorname{int}\Delta_d$. Let $\zeta:\Delta_d\to[0,1]$ be smooth with $\zeta=1$ on $\mathcal O$, put $|D\zeta|=\sup_{\Delta_d}\max_l|(\partial_{m_l}-\partial_{m_d})\zeta|$ and $M_1=2(d-1)M+1$, and assume $T(M_1|D\zeta|+1)<1$. Let $\mathcal K\subset\operatorname{int}\Delta_d$ be compact with $\zeta=0$ on $\mathcal K$, choose $\eta>0$ with $\mathcal K\subset D_\eta$, and put $C=(d-1)M+1$, $\delta'=\frac12\eta e^{-(C+1)T}$ and $a'_\eps=\sup_{[\delta',1]}\phi_\eps$. Then there is $c>0$ such that
\begin{equation}\label{localclassical}
 \sup_{[0,T]\times\mathcal K}|U^{\eps,\phi}-U^{\rm cl}|=O\big(\eps+a'_\eps+e^{-c/\eps}\big).
\end{equation}
\end{lemma}
\begin{proof}
Put $N=D_{\delta'}\setminus\mathcal O$, a compact subset of the interior near which $U^{\rm cl}$ is defined, so $U^{\rm cl}$ and its derivatives of the orders above are bounded on $[0,T]\times N$; also $\mathcal K\subset N$. Let $M_0=M+\sup_{[0,T]\times N}|U^{\rm cl}|$, and take $\eps$ so small that $(d-1)a'_\eps+2\eps\le1$, $2da'_\eps+4\eps\le1$ and $\eps\le\delta'/2$. Write $U=U^{\eps,\phi}$, $E=U-U^{\rm cl}$, $b^{(i)}=b(m,U)+R^\phi+2\eps(e_i-m)$ and $\mathcal P_i=\partial_t+\eps\operatorname{tr}(AD^2\,\cdot\,)+b^{(i)}\cdot D$. Let $\vartheta_{ij}\in[0,1]$ be the difference quotient of $x\mapsto x_+$ between $U_i-U_j$ and $U^{\rm cl}_i-U^{\rm cl}_j$ (say $\vartheta_{ij}=1$ where these coincide), so that $(U_i-U_j)_+-(U^{\rm cl}_i-U^{\rm cl}_j)_+=\vartheta_{ij}(E_i-E_j)$, and similarly for the squares with the factor $(U_i-U_j)_++(U^{\rm cl}_i-U^{\rm cl}_j)_+$. Subtracting the two equations gives
\begin{equation}\label{errorsystem}
 \mathcal P_iE_i+\sum_{j\ne i}(c_{ij}+\phi_\eps(m_j))(E_j-E_i)+(b(U)-b(U^{\rm cl}))\cdot DU^{\rm cl}_i+\rho_i=0,\qquad E(T)=0,
\end{equation}
with $c_{ij}=\frac12\{(U_i-U_j)_++(U^{\rm cl}_i-U^{\rm cl}_j)_+\}\vartheta_{ij}\ge0$ and
\[
 \rho_i=\eps\operatorname{tr}(AD^2U^{\rm cl}_i)+(R^\phi+2\eps(e_i-m))\cdot DU^{\rm cl}_i+\sum_{j\ne i}\phi_\eps(m_j)(U^{\rm cl}_j-U^{\rm cl}_i).
\]
On $N$ every $\phi_\eps(m_j)\le a'_\eps$, so $|R^\phi|_1\le2da'_\eps$ and $\bar\rho=\sup_{[0,T]\times N}\max_i|\rho_i|=O(\eps+a'_\eps)$. Also $b_l(U)-b_l(U^{\rm cl})=\sum_{j\ne l}(m_j\vartheta_{jl}+m_l\vartheta_{lj})(E_j-E_l)$, so $|b(U)-b(U^{\rm cl})|_1\le4(d-1)\max_j|E_j|$. Since $b(U)-b(U^{\rm cl})$ is tangent, the coupling term in \eqref{errorsystem}, which may have either sign, is at most $D_0\max_j|E_j|$ in absolute value, where $D_0=4(d-1)\sup_N\max_{i,l}|(\partial_{m_l}-\partial_{m_d})U^{\rm cl}_i|$. By Lemma~\ref{localization}, $0\le U_i\le M$. Hence $b_l(m,U)\ge-(d-1)Mm_l$, and $\sum_lb_l=0$ gives $|b(m,U)|_1\le2(d-1)M$, so $|b^{(i)}|_1\le M_1$; as in the proof of Lemma~\ref{localization}, $b^{(i)}_j\ge-Cm_j$ on $N$.

Put $C_\zeta=1+\sup_{\Delta_d}\{|\operatorname{tr}(AD^2\zeta)|+D\zeta\cdot AD\zeta\}$, $\kappa_1=\delta'/(4\eps)$, $\kappa_2=1/(2\eps C_\zeta)$, and
\[
 \psi_1=\sum_j\psi_{1,j},\quad \psi_{1,j}=\Big(\frac{\delta'}{m_j}\Big)^{\kappa_1}e^{\kappa_1(C+1)(T-t)},\qquad
 \psi_2=\exp\big(\kappa_2\{\zeta-1+(M_1|D\zeta|+1)(T-t)\}\big).
\]
Since $A_{jj}=m_j(1-m_j)$, $b^{(i)}_j\ge-Cm_j$ and $m_j\ge\delta'$ on $N$,
\[
 \mathcal P_i\psi_{1,j}=\kappa_1\psi_{1,j}\Big\{-(C+1)+\eps(\kappa_1+1)\frac{1-m_j}{m_j}-\frac{b^{(i)}_j}{m_j}\Big\}
 \le\kappa_1\psi_{1,j}\Big\{-\frac34+\frac\eps{\delta'}\Big\}\le0.
\]
Since $b^{(i)}$ is tangent, $|b^{(i)}\cdot D\zeta|\le|b^{(i)}|_1|D\zeta|\le M_1|D\zeta|$, and therefore
\[
 \mathcal P_i\psi_2=\kappa_2\psi_2\big\{-M_1|D\zeta|-1+\eps\operatorname{tr}(AD^2\zeta)+\eps\kappa_2D\zeta\cdot AD\zeta+b^{(i)}\cdot D\zeta\big\}
 \le\kappa_2\psi_2\big\{-\tfrac12+\eps C_\zeta\big\}\le0
\]
once $\eps C_\zeta\le1/2$. Moreover $\psi_1\ge1$ where some $m_j=\delta'$, $\psi_2\ge1$ on the closure of $\mathcal O$, and on $\mathcal K$, where all $m_j\ge\eta$ and $\zeta=0$,
\[
 \psi_1+\psi_2\le d\,2^{-\kappa_1}+e^{-\kappa_2\{1-T(M_1|D\zeta|+1)\}}.
\]

Fix $\Lambda\ge D_0$ and $\varpi>0$, and put $\Gamma=\bar\rho(T-t)+M_0(\psi_1+\psi_2)+\varpi(1+T-t)$. We claim that $\pm E_ie^{-\Lambda(T-t)}\le\Gamma$ on $[0,T]\times N$ for every $i$ and both signs. Otherwise the maximum of $\pm E_ie^{-\Lambda(T-t)}-\Gamma$ over $t$, $m\in N$, $i$ and the sign is positive; let it be attained at $(t_*,m_*)$, for the index $i$ and the sign $\pm$, and put $e=\pm E_ie^{-\Lambda(T-t)}$. Since $E(T)=0$, and $|E|\le M_0\le M_0(\psi_1+\psi_2)$ at points of $N$ where some $m_j=\delta'$ or which lie in the closure of $\mathcal O$, we have $t_*<T$, and $m_*$ is an interior point of $N$. At $(t_*,m_*)$ we have $\partial_t(e-\Gamma)\le0$, $D(e-\Gamma)=0$ and $D^2(e-\Gamma)\le0$ on tangent vectors, so $\mathcal P_ie\le\mathcal P_i\Gamma$. There also $e>0$ and $e\ge e^{-\Lambda(T-t)}|E_j|$ for every $j$, so the switching sum in \eqref{errorsystem}, multiplied by $\pm e^{-\Lambda(T-t)}$, is nonpositive, and \eqref{errorsystem} gives
\[
 \mathcal P_ie=\Lambda e\mp e^{-\Lambda(T-t)}\Big[\sum_{j\ne i}(c_{ij}+\phi_\eps(m_j))(E_j-E_i)+(b(U)-b(U^{\rm cl}))\cdot DU^{\rm cl}_i+\rho_i\Big]\ge(\Lambda-D_0)e-\bar\rho\ge-\bar\rho.
\]
But $\mathcal P_i\Gamma=-\bar\rho+M_0\mathcal P_i(\psi_1+\psi_2)-\varpi\le-\bar\rho-\varpi$, a contradiction. Letting $\varpi\downarrow0$ gives $|E|\le e^{\Lambda T}\{\bar\rho T+M_0(\psi_1+\psi_2)\}$ on $[0,T]\times N$, which proves \eqref{localclassical} with $c=\min\{\delta'\log2/4,\ (1-T(M_1|D\zeta|+1))/(2C_\zeta)\}$. No bound on derivatives of $U^{\eps,\phi}$ uniform in $\eps$ was used.
\end{proof}

\section{Reduction and the regular collars}
Write $u=U_1$. On an interior compact cylinder, \eqref{passive} makes all passive components tend uniformly to zero and leaves $u$ uniformly positive. Elementary algebra in the rates and Hamiltonians gives
\begin{equation}\label{scalarperturb}
 u_s+q\partial_p(pu^2)
 =\eps\mathcal L u+E_\eps\cdot D_{p,r}u+F_\eps,
 \qquad \|E_\eps\|_\infty+\|F_\eps\|_\infty\longrightarrow0.
\end{equation}
Here $\mathcal L$ is the scalar Kimura operator including the extra drift $2(e_1-m)\cdot D$, and the errors contain no derivatives of passive values. Indeed,
$b_1=-2qpu+O(\max_{i>1}U_i)$, $b_j-b_d=O(\max_{i>1}U_i)$, $H_1=qu^2+O(\max_{i>1}U_i)$, and the repulsion terms have uniformly small coefficients on the compact set. Treat these error functions, obtained from the actual solution, as fixed coefficients in \eqref{scalarperturb}.

Under $V=qyu$, the unperturbed part of \eqref{scalarperturb} becomes $V_s+VV_y$. In $(y,r)$ the second-order coefficients are
\begin{equation}\label{coefficients}
 A_{yy}=\frac{1-y^2}{4},\qquad A_{yr_j}=-\frac{yr_j}{2},\qquad
 A_{r_jr_l}=m_j\delta_{jl}+m_d-r_jr_l.
\end{equation}
The first-order coefficients are $(1-y^2)/(4y)$ and $-r_j$, and the zeroth-order coefficient is $-(1-y^2)/(4y^2)$. After the fixed change $\theta=y-1/2-\alpha r_2$, \eqref{scalarperturb} is exactly
\begin{equation}\label{Veq}
 V_s+VV_\theta=\eps\cD V+B_\eps\cdot D_{\theta,r}V+C_\eps,
 \qquad \|B_\eps\|_\infty+\|C_\eps\|_\infty\longrightarrow0.
\end{equation}
On interior compact sets, $\cD$ has fixed smooth coefficients and a positive definite second-order matrix. Denote its $\partial_{\theta\theta}$ coefficient by $a(\theta,r)>0$. The conversion of the errors is harmless because $y$ is bounded away from zero.

Consider the rectangles
\[
 \Omega=\{(\theta,r):|\theta|<1/8,\ |r|_\infty<R_d\},\qquad
 \Omega^\pm=\{(\theta,r):|\theta|<1/8\pm1/64,\ |r|_\infty<(1\pm1/8)R_d\},
\]
and the collar $\mathcal C=\overline{\Omega^+}\setminus\Omega^-$, a compact neighborhood of $\partial\Omega$. On $\overline{\Omega^+}$ one has $11/32<y<2/3$ and $L(r)^2\ge1-(2d-3)\cdot9R_d/8>81/100$, so $\overline{\Omega^+}$ lies in the simplex interior. Since $\chi=1$ wherever $y\ge1/8$, the initial value for \eqref{Veq} is exactly $2-\beta(r)\tanh(k\theta')$ at $(\theta',r)$ whenever $(\theta,r)\in\overline{\Omega^+}$ and $\theta-1/32\le\theta'\le\theta$.

We now choose $T_0$, before choosing $k$. All derivatives through any fixed order of the initial data are bounded uniformly in $k$ outside the smaller strip
\[
 \mathcal S=\{|\theta|\le1/16,\ |r|_\infty\le R_d/2\}.
\]
This follows from $\sup_{k\ge1}k^j\operatorname{sech}^2(k/16)<\infty$; in the remaining region $\beta$ vanishes. The cutoff near $p=0$ is fixed and stays away from this strip. Put $\mathcal S_1=\{|\theta|\le3/32,\ |r|_\infty\le3R_d/4\}$. Then $\mathcal S\subset\mathcal S_1\subset\Omega^-$, so $\mathcal S_1$ does not meet $\mathcal C$. Fix a smooth $\zeta:\Delta_d\to[0,1]$ with $\zeta=1$ on $\mathcal S_1$ and $\zeta=0$ outside $\Omega^-$, for instance a product of cutoff functions of $\theta$ and of $r_2,\ldots,r_{d-1}$; it depends only on $d$.

Let $C_1\ge\sup_{k\ge1}k\operatorname{sech}^2(k/16)$ bound the negative part of $\partial_yv_0$ outside $\mathcal S$, uniformly in $k$. With $M=34/q$, and with $M_1$ and $|D\zeta|$ as in Lemma~\ref{consistency}, choose $T_0>0$ with
\[
 T_0C_1<1/2,\qquad 3T_0<1/32,\qquad T_0(M_1|D\zeta|+1)<1.
\]
These are finitely many positive requirements independent of $k$.

For the deterministic ansatz $(u,0,\ldots,0)$ the equations reduce exactly to \eqref{burgers}, with $r$ constant along characteristics. Its characteristic map is $y=z+sv_0(z,r)$. Let $0<T\le T_0$ and $0<s\le T$. Since $0\le v_0\le3$, every Hopf--Lax minimizer $z$ at $(s,y,r)$ satisfies $y-z=sv_0(z,r)\in[0,3s]$. The set $\mathcal S'=\{-1/16\le\theta\le1/16+3T,\ |r|_\infty\le R_d/2\}$ lies in the interior of $\mathcal S_1$, since $3T<1/32$. If $(\theta,r)\notin\mathcal S'$, the footpoint segment $\{(\theta',r):\theta-3s\le\theta'\le\theta\}$ does not meet $\mathcal S$: this is clear if $|r|_\infty>R_d/2$, and the segment lies to the left of $\mathcal S$ if $\theta<-1/16$ and to its right if $\theta>1/16+3T$. On this segment the derivative of the action $z\mapsto\int_0^zv_0+(y-z)^2/(2s)$ increases at rate at least $1/s-C_1\ge1/(2s)$, so the minimizer is unique and nondegenerate, and $\partial_z(z+sv_0)\ge1-sC_1\ge1/2$. By the implicit function theorem $v$ is smooth near each such point, up to $s=0$. Hence $U^*$ is a classical solution of the deterministic master equation with terminal value $g$ on $[0,T]\times(\operatorname{int}\Delta_d\setminus\mathcal S')$, a neighborhood of $[0,T]\times(\operatorname{int}\Delta_d\setminus\operatorname{int}\mathcal S_1)$.

Apply Lemma~\ref{consistency} with $\mathcal O=\operatorname{int}\mathcal S_1$, $U^{\rm cl}=U^*$, the function $\zeta$, and any compact $\mathcal K\subset\operatorname{int}\Delta_d\setminus\Omega^-$; its hypotheses hold by the choice of $T_0$. Thus $U^{\eps,\phi}\to U^*$ uniformly on $[0,T]\times\mathcal K$. In particular the actual boundary values on the collar $\mathcal C$ converge uniformly to the classical field. No classical field is assumed across the shock.

\section{A uniform Cole--Hopf estimate}
Let $w(s,x;\nu,\lambda)$ solve viscous Burgers on $\R$, with viscosity $\nu>0$ and initial value $2-\lambda\tanh(kx)$, where $0\le\lambda\le1$. Write
\begin{equation}\label{action}
 S(z)=2z-\frac\lambda k\log\cosh(kz)+\frac{(x-z)^2}{2s},\qquad
 w=\frac{x-\E_\nu z}{s},\qquad
 d\mu_\nu(z)=\frac{e^{-S(z)/(2\nu)}dz}{\int e^{-S/(2\nu)}}.
\end{equation}
This is the Cole--Hopf formula; differentiating the heat kernel verifies the equation and initial value. In this section expectations and covariances use $\mu_\nu$. All estimates are uniform for $x$ in a fixed compact interval, $s$ in a fixed bounded interval, and $\lambda\in[0,1]$. The parameter $k$ is fixed.

\begin{lemma}\label{CH}
Put $Q=-w_x\ge0$. There is a modulus $\omega(\nu)\to0$ such that
\begin{align}
 |w|&\le3,\qquad |\nu w_{xx}|+|w_\lambda|\le C(1+Q),\label{CH1}\\
 |\nu w_\nu|+|\nu^2w_{x\nu}|+|\nu^3w_{\nu\nu}|
 +|\nu^2w_{\nu\lambda}|+|\nu w_{x\lambda}|+|\nu w_{\lambda\lambda}|
 &\le\omega(\nu)(1+Q).\label{CH2}
\end{align}
The estimates hold from $s=0$ by continuous extension.
\end{lemma}
\begin{proof}
The maximum principle gives $1\le w\le3$ and $w_x\le0$. For $s\ge s_0>0$, put $Z=z-\E z$, $A=S-\E S$, and $R=\partial_\lambda S=-k^{-1}\log\cosh(kz)$. Differentiation under the integral gives
\begin{align*}
 \Var(z)&=2\nu s^2(Q+1/s),&
 w_{xx}&=-\frac{\E Z^3}{4\nu^2s^3},\\
 w_\nu&=-\frac{\E ZA}{2\nu^2s},&
 w_{x\nu}&=\frac{\Var(z)}{2\nu^2s^2}-\frac{\E Z^2A}{4\nu^3s^2},\\
 w_{\nu\nu}&=\frac{\E ZA}{\nu^3s}-\frac{\E ZA^2}{4\nu^4s},&
 w_\lambda&=\frac{\Cov(z,R)}{2\nu s},\\
 w_{x\lambda}&=\frac{\E Z^2(R-\E R)}{4\nu^2s^2},&
 w_{\lambda\lambda}&=-\frac{\E Z(R-\E R)^2}{4\nu^2s},\\
 w_{\nu\lambda}&=-\frac{\Cov(z,R)}{2\nu^2s}
 +\frac{\E Z(R-\E R)A}{4\nu^3s}.&&
\end{align*}

The moment estimates below are uniform, including the possibly degenerate transition at shock formation. On a bounded $z$ interval the derivatives $S'',S''',S''''$ do not vanish simultaneously: if $\lambda=0$ then $S''=1/s$; otherwise $S'''$ vanishes only at zero, where $S''''=2\lambda k^3>0$. Outside a sufficiently large interval $S''$ is uniformly positive. A finite cover consequently gives
\[
 |\{z:S(z)-\min S\le h\}|\le C h^{1/4}\quad(0<h\le1).
\]
The elementary sublevel estimate used here follows by choosing $j+1$ separated points of a set on which $|F|\le h$. Divided differences bound $\inf|F^{(j)}|$ by $C_jh/|\{|F|\le h\}|^j$. Apply this on each covering interval where one derivative of order $j\le4$ has a fixed nonzero sign. The quadratic tails of $S$ handle larger sublevels.

At a global minimizer, the uniform upper bound on $S''$ gives a Gaussian lower bound $c\sqrt\nu$ for the partition integral after subtraction of $\min S$. Integrating the sublevel estimate, or summing energy strips of width $\nu$, gives
\begin{equation}\label{actionmom}
\E(S-\min S)^j+\E|A|^j\le C\nu^{j-1/4}\qquad(1\le j\le4).
\end{equation}
Moreover $\E z\in[x-3s,x-s]$, by integration by parts in \eqref{action}. Gaussian far tails, together with $\Var(z)\ge2\nu s$, imply
\begin{equation}\label{zmom}
 \E|Z|^3+\E Z^4\le C\Var(z).
\end{equation}
Indeed, on a fixed bounded interval the higher powers are bounded by a constant times $Z^2$; the complementary moments are exponentially small compared with $\nu$. Equations \eqref{actionmom}--\eqref{zmom} and Cauchy--Schwarz prove the viscosity-derivative bounds in \eqref{CH2}, for example with factor $C\nu^{3/8}$, and the bound on $\nu w_{xx}$. Since $R$ is $1$-Lipschitz, $|\Cov(z,R)|\le\Var(z)$ and $\E[Z^2(R-\E R)^2]\le C\Var(z)$. This also proves the bounds on $w_\lambda$ and $\nu^2w_{\nu\lambda}$.

The remaining two derivatives need a sharper observation. Set $h=x-2s$. Apart from an additive constant,
\[
 S(z)=z^2/(2s)-hz/s-(\lambda/k)\log\cosh(kz).
\]
If $h\ne0$, there is a unique global minimizer, of the sign of $h$. To see this for $h>0$, compare $z$ with $-z$ and then use that $S''$ is increasing on the positive half-line and $S'(0)<0$. If $h=0$, the minimizer is zero or a symmetric pair. Hence the even function $R$ has the same value at every global minimizer, for every parameter.

Compactness now implies that the oscillation of $R$ on the sublevel $S-\min S\le a$ tends uniformly to zero as $a\downarrow0$. Outside a fixed such sublevel the Gibbs moments are exponentially small. On the sublevel use this small oscillation for one factor and the Lipschitz covariance bound for the other factors. Since $\Var(z)\ge c\nu$, it follows that
\[
 \E[Z^2(R-\E R)]=o(\Var(z)),\qquad
 \E[Z(R-\E R)^2]=o(\Var(z)),
\]
uniformly in the parameters. For the second estimate one uses
$\E|Z||R-\E R|\le\Var(z)$; centering by $\E R$ changes the sublevel oscillation only by the same vanishing modulus and an exponentially small tail. The displayed derivative identities prove the desired bounds on $\nu w_{x\lambda}$ and $\nu w_{\lambda\lambda}$.

Finally choose $s_0\le1/(2k)$. Before this time the differentiated Burgers equation gives $-2k\le w_x\le0$. Successive spatial derivatives satisfy linear parabolic equations with bounded coefficients and already bounded lower derivatives. Differentiating in $\nu$ or $\lambda$ gives the same type of equations, with initial data zero or derivatives of $2-\lambda\tanh(kx)$; the source for $w_\nu$, for example, is $w_{xx}$. Induction through the spatial order four and parameter order two gives bounds independent of $\nu$ on this initial interval. Each scaled derivative in \eqref{CH2} then tends to zero there as well. This completes the estimates through $s=0$.
\end{proof}

\section{Comparison through the shock}
On $\Omega^+$, set
\[
 W_\eps(s,\theta,r)=w(s,\theta;\eps a(\theta,r),\beta(r)).
\]
The chain rule and Lemma~\ref{CH} give, with a modulus $\eta_\eps\to0$,
\begin{align}
 1-W_\theta&\asymp1+Q,\qquad W_\theta\le\eta_\eps,
 \qquad |DW|+\eps|D^2W|\le C(1-W_\theta),\label{weighted}\\
 |W_s+WW_\theta-\eps\cD W|&\le\eta_\eps(1-W_\theta).\label{residual}
\end{align}
The potentially large derivatives cause no omitted term. Since $\beta$ is independent of $\theta$, $W_\theta=w_x+O(\nu w_\nu)$ with $\nu=\eps a$. The main normal diffusion cancels $w_s+ww_x=\nu w_{xx}$. The remaining normal terms involve $\nu w_\nu$, $\nu^2w_{x\nu}$ and $\nu^3w_{\nu\nu}$. Mixed and tangential diffusion adds $\nu w_{x\lambda}$, $\nu w_{\lambda\lambda}$ and $\nu^2w_{\nu\lambda}$, all small in the weighted sense. Terms with a single $w_\lambda$ carry a factor $\eps$. The fixed lower-order coefficients are bounded. This proves \eqref{weighted}--\eqref{residual}.

On $[0,T]\times\mathcal C$ the reference converges uniformly to the classical field $v$. If $|r|_\infty\ge7R_d/8$, then $\beta(r)=0$ and $W_\eps=2=v$. Otherwise $|\theta|\ge7/64$; take $x=\theta$, $\lambda=\beta(r)$ and $0<s\le T$ in \eqref{action}. Since $1\le2-\lambda\tanh(kz)\le3$, we have $(z-x+s)/s\le S'(z)\le(z-x+3s)/s$, so $S$ decreases on $(-\infty,x-3s]$ and increases on $[x-s,\infty)$, while on $[x-3s,x-s]$ one has $|z|>1/16$, because $3s<1/32$, and hence $S''\ge1/s-C_1\ge1/(2s)$. Thus $S$ has a unique minimizer $z_*$, the classical value is $(x-z_*)/s=v(s,\theta,r)$, the data agreeing on $[\theta-3s,\theta]$ by the choice of $\Omega^+$, and integrating these bounds gives $S(z)-S(z_*)\ge(z-z_*)^2/(6s)$ for all $z$. Also $S''\le1/s$ and $|S'''|\le k^2$. Pairing the Gibbs weights at $z_*\pm\xi$ and using $|e^{-a}-e^{-b}|\le|a-b|e^{-\min(a,b)}$ and Taylor's formula, we get
\[
 |\E_\nu z-z_*|\le\frac{\int_0^\infty\xi\cdot\frac{k^2\xi^3}{6\nu}\,e^{-\xi^2/(12\nu s)}\,d\xi}{\int_{\R}e^{-\xi^2/(4\nu s)}\,d\xi}\le16k^2\nu s^2,
\]
so $|W_\eps-v|\le16k^2\eps s\,a(\theta,r)$. Let $b_\eps\to0$ be the sum of the uniform errors of $W_\eps$ and of $V_\eps$ with respect to $v$ on $[0,T]\times\mathcal C$, the second tending to zero by Section~3, and let $L$ bound $|\partial_\theta v|$ there. Let $\Delta_\eps\to0$ bound the two error coefficients in \eqref{Veq} on the rectangle.

Choose nonnegative functions $a_0(s),h(s)$ by
\begin{equation}\label{odebarrier}
 K=C(h+a_0+\eta_\eps+\Delta_\eps),\quad
 h'=a_0+K,\quad a_0'=Lh'+K,\quad
 h(0)=0,\quad a_0(0)=b_\eps,
\end{equation}
where $C$ is a sufficiently large fixed constant. They tend uniformly to zero. For small $\eps$ we have $h<1/64$ on $[0,T]$; then $(\theta\mp h(s),r)\in\Omega^+$ for $(\theta,r)\in\overline\Omega$, and for $(\theta,r)\in\partial\Omega$ the segment from $(\theta,r)$ to $(\theta\mp h(s),r)$ lies in $\mathcal C$. The functions
\[
 W^+(s,\theta,r)=W_\eps(s,\theta-h(s),r)+a_0(s),\qquad
 W^-(s,\theta,r)=W_\eps(s,\theta+h(s),r)-a_0(s)
\]
are respectively a supersolution and a subsolution of \eqref{Veq}. Indeed, translation of the fixed smooth coefficients changes the residual by at most $Ch(1-W_\theta)$, using the Hessian bound in \eqref{weighted}. The other errors, including the zeroth-order term on $a_0$, are at most $K(1-W_\theta)$. The upper residual is therefore at least
\[
 (-h'+a_0)W_\theta+a_0'-K(1-W_\theta)=a_0'-K\ge0;
\]
the lower residual is at most $K-a_0'\le0$.

The initial inequalities hold because $W_\eps(0)=V_\eps(0)$ on the rectangle. On its boundary, \eqref{odebarrier} gives $a_0-Lh\ge b_\eps$, which proves the boundary inequalities. At each fixed $\eps$ the scalar equation is uniformly parabolic on this interior rectangle, so the comparison principle yields $W^-\le V_\eps\le W^+$. Only boundedness, not derivatives, of the error coefficients was used.

For each $r$, the reference tends to the Burgers entropy solution, except possibly on the shock plane $\theta=2s$ when $\beta(r)ks>1$. This is also the solution defined from the full data \eqref{data}: both sets of minimizing footpoints lie within distance $3T_0$ of the observation point, where the initial functions coincide; their primitives differ only by an additive constant. Boundedness gives convergence in local $L^1$. Also $W_\theta\le\eta_\eps$ and $1\le W\le3$ give a uniform bound on its total variation on each fixed $\theta$ interval. Its small normal shifts therefore tend to zero in $L^1$, uniformly in the other variables. The comparison just proved gives $V_\eps\to v$ in $L^1$ on the rectangle. Outside it, Section~3 gives local uniform convergence, because for every compact $\mathcal K\subset\operatorname{int}\Delta_d$ the closure of $\mathcal K\setminus\Omega$ is a compact subset of $\operatorname{int}\Delta_d\setminus\Omega^-$. Lemma~\ref{localization} removes the passive components. This proves \eqref{convergence} for every admissible family.

\section{Admissibility, nonemptiness and scope}
The change $(p,r)\mapsto(y,r)$ has a smooth nonzero Jacobian on interior compact sets, so the convergence above is in the simplex measure. It implies compactness and pairwise independence of the regularizations. The entropy condition \eqref{burgers}, imposed for almost every parameter $r$, is intrinsic to the limiting deterministic equation and has no reference to $\eps$ or $\phi$. All selected feedback fields are bounded by the value estimate.

Entropy uniqueness also applies on the parameter-dependent interval used here. The local Kato inequality for two nonnegative entropy solutions has flux
$\tfrac12(v+\widetilde v)|v-\widetilde v|\ge0$. Integrate with a cutoff that has compact support in $(0,L(r))$, increases near zero and decreases near $L(r)$. Bounded feedback gives $v,\widetilde v=O(y)$ near zero, so the left cutoff error vanishes. The right cutoff contribution is nonpositive and can be discarded. The initial $L^1$ distance is zero; hence the solutions agree. This is the usual entropy contraction argument \cite[Theorem~1]{Kru}, with the inflow and outflow terms made explicit. It proves uniqueness under the stated intrinsic condition.

If a classical full master solution exists, follow one of its characteristics. The value vector solves $\dot u_i=H_i(u)$ backwards from $(G_k,0,\ldots,0)$. Uniqueness of this locally Lipschitz ODE keeps the passive values zero and the active value positive. The full classical solution therefore has the scalar reduction \eqref{burgers}, and agrees with its entropy solution. This verifies classical consistency without assuming the ansatz for an arbitrary candidate classical field.

The admissible class is nonempty. Bayraktar, Cecchin, Cohen and Delarue \cite[Theorem~3.4, equation~(2.19)]{BCCD} provide a classical Wright--Fisher solution for smooth costs when the repulsion is sufficiently large near zero and vanishes beyond an arbitrarily small boundary layer. Their condition (2.19) asks only that the repulsion be nonincreasing and Lipschitz, equal to a large constant on $[0,\delta]$ and to zero on $(2\delta,\infty)$, so it may be taken smooth. Their noise parameter is $\sqrt{2\eps}$ in our normalization. Their first derivatives extend continuously to the boundary.

Uniqueness extends to the class used here. Let $U$ be this solution, let $\widetilde U$ be any solution of \eqref{master} with the regularity required in the definition of $\cA$, and put $E=U-\widetilde U$. The subtraction leading to \eqref{errorsystem}, with $U^{\rm cl}$ replaced by $\widetilde U$ and now on the whole cylinder, gives
\[
 \mathcal P_iE_i+\sum_{j\ne i}(c_{ij}+\phi_\eps(m_j))(E_j-E_i)+(b(U)-b(\widetilde U))\cdot D\widetilde U_i=0,\qquad E(T)=0,
\]
with $c_{ij}\ge0$. The last term may have either sign; it is at most $D_0\max_j|E_j|$ in absolute value, where $D_0$ is $4(d-1)$ times a bound for the first derivatives of $\widetilde U$. Fix $\Lambda\ge D_0$, $0<a<1$, $h>0$ and $\kappa>0$, and maximize $\pm E_ie^{-\Lambda(T-t)}+hB_a(m)-\kappa(T-t)$ over $t$, $m\in\Delta_d$, $i$ and the sign. As in the proof of Lemma~\ref{localization}, bounded first derivatives of $E$ place the maximum in the interior. Suppose that it occurs at some $t<T$, and let $e=\pm E_ie^{-\Lambda(T-t)}$ for the maximizing index and sign. Then $e\ge e^{-\Lambda(T-t)}|E_j|$ for every $j$, so, as in the proof of Lemma~\ref{consistency}, $\mathcal P_ie\ge(\Lambda-D_0)e\ge0$ there. On the other hand, at this maximum the first-order condition and the estimates of the terms of $hB_a$ in the proof of Lemma~\ref{localization} give $\mathcal P_ie\le-\kappa+\eps d(1-a)(L_E+had^{1-a})+O(h)$, where $L_E$ bounds the first derivatives of $E$. For $\kappa$ larger than the last two terms this is impossible. Hence the maximum occurs at $t=T$, where $E=0$, and $\pm E_i\le e^{\Lambda T}(hd^{1-a}+\kappa T)$. Letting $h\downarrow0$ and $\kappa\downarrow\eps d(1-a)L_E$, and then $a\uparrow1$, gives $E=0$. Taking the layer width to zero gives an admissible family. Only sufficiently small $\eps$ matters for the limit; if a family is required at larger $\eps$ too, the substitution $U=cV$, $\tau=ct$, $\phi=c\widetilde\phi$, with $c>2\eps$, reduces it to the same theorem.

At $p=1/4,r=0$, the intrinsic terminal covector is $G_k\,dm_1$. Varying $m_2$ with $m_d$ compensating gives
\[
 \partial_{m_2}G_k=\frac{4\alpha k}{q}\ne0,
 \qquad \partial_{m_1}(g_2-g_d)=0.
\]
It is not a gradient, even after a common component is added to the costs. On the core where $\beta=1$, the initial Burgers derivative is $-k$ at $\theta=0$. Characteristics first cross at $s=1/k$; for later times the symmetric decreasing front has its entropy shock at $\theta=2s$. The construction of $T_0$ did not depend on $k$, so this shock occurs strictly inside the permitted cylinder whenever $kT>1$.

The theorem is a sufficient-condition result for the stated data family, not a claim that arbitrary nonpotential costs always have a unique vanishing-noise limit. No cost correction, prescribed rate for $\phi_\eps$, or uniform bound on noisy derivatives is part of the theorem.

\begin{corollary}[Every finite horizon]\label{allhorizons}
For any $d\ge3$ and $T>0$, choose $\sigma>0$ with $\sigma T\le T_0$ and use $f=0$, $g=(\sigma G_k,0,\ldots,0)$. All the existence, independence, intrinsic admissibility and classical-consistency conclusions of Theorem~\ref{main} hold. The entropy initial value is $\sigma v_0$, and an interior shock occurs when $\sigma kT>1$.
\end{corollary}
\begin{proof}
Set $\widehat U(\tau,m)=\sigma^{-1}U(\tau/\sigma,m)$, $\widehat\eps=\eps/\sigma$, and
$\widehat\phi_{\widehat\eps}=\sigma^{-1}\phi_{\sigma\widehat\eps}$. The homogeneities of $H$ and $b$ show directly that \eqref{master} becomes the same equation on the horizon $\sigma T$, with terminal value $(G_k,0,\ldots,0)$. This transformation is a bijection of the respective admissible regularization classes; their regularity, uniqueness and local vanishing properties are preserved. Theorem~\ref{main} applies, and scalar entropy solutions scale as $v^\sigma(s,y,r)=\sigma v(\sigma s,y,r)$. One may take $\sigma=\min(1,T_0/(2T))$ and then choose $k>1/(\sigma T)$ to obtain a shock on any prescribed horizon.
\end{proof}

\section*{Use of AI tools}
The author used research AI agents that the author built and configured (GPT-6 Astra max,
Claude Opus 5.5 max and GPT-6 Astra Pro) during the proof work, for proof search and
construction, proof checking, translation, Python computation and similar tasks. The final
manuscript was written, and the proofs were checked, by the human author. All responsibility
for this work rests with the author personally.

\begin{thebibliography}{9}
\small
\bibitem{CD} P. Cardaliaguet and F. Delarue.
\emph{Selected topics in mean field games}.
Proceedings of the International Congress of Mathematicians 2022, vol.~5, 3660--3703.
DOI: \href{https://doi.org/10.4171/ICM2022/17}{10.4171/ICM2022/17}.
\bibitem{BCCD} E. Bayraktar, A. Cecchin, A. Cohen and F. Delarue.
\emph{Finite state mean field games with Wright--Fisher common noise}.
\href{https://arxiv.org/pdf/1912.06701v2}{arXiv:1912.06701v2}, 10 January 2021, Theorem~3.4 and \S3.2.2.
\bibitem{Kru} S. N. Kruzhkov.
\emph{First order quasilinear equations in several independent variables}.
Mathematics of the USSR-Sbornik 10(2) (1970), 217--243.
DOI: \href{https://doi.org/10.1070/SM1970v010n02ABEH002156}{10.1070/SM1970v010n02ABEH002156}.
\end{thebibliography}
\end{document}
